% Chapter 19: The Sensitivity Conjecture
\let\LocalMacro\relax

\chapter{The Sensitivity Conjecture}

\section{The Problem}
For a boolean function $f: \{0,1\}^n \to \{0,1\}$, several complexity measures characterize how ``complex'' $f$ is:
\begin{itemize}
\item \textbf{Sensitivity} $s(f)$: the maximum number of coordinates whose flip changes $f$.
\item \textbf{Block sensitivity} $bs(f)$: the maximum number of disjoint blocks of coordinates whose simultaneous flip changes $f$.
\item \textbf{Degree} $\deg(f)$: the degree of the real polynomial representing $f$.
\end{itemize}

It was known that $\deg(f) \le s(f)^2$ and $bs(f) \le \deg(f)^2$, so the conjecture would follow from $s(f) \le \deg(f)^c$.

\begin{conjecture}[Nisan-Szegedy, 1992]
$s(f)$ and $bs(f)$ are polynomially related: $bs(f) \le s(f)^c$ for some constant $c$.
\end{conjecture}

\begin{historicalbox}
The sensitivity conjecture stood for 27 years, resisting attacks from Fourier analysis, probability theory, and complexity theory. It was one of the oldest open problems in theoretical computer science.
\end{historicalbox}


\section{Mathematical Theory}

\subsection{Prerequisites}
This chapter assumes familiarity with:
\begin{itemize}
\item Boolean functions ($f: \{0,1\}^n \to \{0,1\}$)
\item Spectral graph theory (eigenvalues of adjacency matrices)
\item Basic complexity theory (decision trees, polynomial degree)
\end{itemize}

\subsection{The Discovery}

\subsection{Huang (2019)}

\begin{theorem}[Huang \cite{huang2019}]
$s(f) \ge \sqrt{\deg(f)}$, hence $bs(f) \le s(f)^4$.
\end{theorem}

Huang's proof is famously short (2 pages) and uses only spectral graph theory:

\begin{enumerate}
\item Consider the $n$-dimensional hypercube graph $Q_n$, whose vertices are $\{0,1\}^n$ and edges connect strings differing in one coordinate.
\item Construct a signing of the adjacency matrix $A$ of $Q_n$ such that $A^2 = nI$ (a Hadamard matrix structure).
\item By the Cauchy interlacing theorem, the subgraph induced by $f^{-1}(1)$ has an eigenvalue $\ge \sqrt{n}$.
\item This implies the subgraph has minimum degree $\ge \sqrt{n}$, so $s(f) \ge \sqrt{n}$.
\item Since $\deg(f) = n$ for the dictator function, $s(f) \ge \sqrt{\deg(f)}$.
\end{enumerate}

\begin{highlightbox}
Huang's proof resolved a 27-year-old problem using a single clever idea: a signing of the hypercube's adjacency matrix. The proof uses no Fourier analysis, no probability, and no complexity theory---just linear algebra.
\end{highlightbox}

\subsection{Impact}
\begin{itemize}
\item \textbf{Complexity theory}: Established polynomial relationships between all known boolean function complexity measures.
\item \textbf{Spectral methods}: The signing technique has been applied to other graph-theoretic problems.
\item \textbf{Open question}: The exact polynomial relating $s(f)$ and $bs(f)$ remains unknown---Huang's proof gives $c = 4$, but the conjectured value is $c = 1$.
\end{itemize}

\subsection{Exercises}

\begin{exercise}[Basic]
Compute $s(f)$ and $bs(f)$ for the parity function $f(x_1, \ldots, x_n) = x_1 \oplus \cdots \oplus x_n$.
\end{exercise}

\begin{exercise}[Intermediate]
Explain the Cauchy interlacing theorem: if $A$ is an $n \times n$ symmetric matrix and $B$ is an $(n-1) \times (n-1)$ principal submatrix, how do their eigenvalues relate?
\end{exercise}

\begin{exercise}[Challenge]
Sketch Huang's construction of the signing of $Q_n$'s adjacency matrix. Why does $A^2 = nI$ hold?
\end{exercise}

\subsection{Timeline}

\begin{tabularx}{\linewidth}{lX}
\toprule
\textbf{Year} & \textbf{Event} \\ \midrule
1992 & Nisan-Szegedy formulate the sensitivity conjecture \\
2019 & Huang proves $s(f) \ge \sqrt{\deg(f)}$ in 2 pages \\
\bottomrule
\end{tabularx}

\subsection{Citations}

\begin{enumerate}
\item Huang, H. (2019). ``Induced subgraphs of hypercubes and a proof of the sensitivity conjecture.'' Annals of Mathematics, 190(3), 949--955.
\item Nisan, N., \& Szegedy, M. (1992). ``On the degree of a boolean function as a real polynomial.'' Computational Complexity, 4(4), 301--313.
\end{enumerate}
